% !TEX root = ../algebra_02.tex
\section{Лемма Бёрнсайда}

\begin{Definition}{Неподвижные точки}{}
	Пусть $G \curvearrowright M$.
	Для элемента $g \in G$ множество его неподвижных точек обозначается как:
	\[ M^g := \{m \in M \mid gm = m\} \] 
\end{Definition}

\begin{Lemma}{Лемма Бёрнсайда}{}
	Пусть $G$ и $M$ конечны.
	Тогда число орбит $N$ действия $G$ на $M$ равно среднему числу неподвижных точек:
	\[ N = \frac{1}{|G|} \sum_{g \in G} |M^g| \] 
\end{Lemma}

\begin{proof}
	Рассмотрим множество пар $S = \{(g, m) \in G \times M \mid gm = m\}$. Посчитаем его мощность двумя способами.
	С одной стороны, суммируя по элементам группы:
	\[ |S| = \sum_{g \in G} |M^g| \] 
	С другой стороны, суммируя по элементам множества:
	\[ |S| = \sum_{m \in M} |\{g \in G \mid gm = m\}| = \sum_{m \in M} |St_m| \] 
	По теореме об орбите и стабилизаторе $|St_m| = \frac{|G|}{|Gm|}$.
	Подставим это в сумму:
	\[ \sum_{m \in M} \frac{|G|}{|Gm|} = |G| \sum_{m \in M} \frac{1}{|Gm|} \] 
	Сгруппируем слагаемые по орбитам.
	Пусть $X = \{Gm \mid m \in M\}$ — множество всех орбит.
	\[ |G| \sum_{O \in X} \sum_{m \in O} \frac{1}{|O|} = |G| \sum_{O \in X} \left( |O| \cdot \frac{1}{|O|} \right) = |G| \sum_{O \in X} 1 = |G| \cdot |X| = |G| \cdot N \] 
	Приравнивая два способа подсчета $|S|$, получаем:
	\[ \sum_{g \in G} |M^g| = |G| \cdot N \implies N = \frac{1}{|G|} \sum_{g \in G} |M^g| \] 
\end{proof}

\begin{Example}{Раскраска ожерелья (группа $D_9$)}{}
	Рассмотрим задачу о количестве способов раскрасить правильный 9-угольник (ожерелье из 9 бусин) в 2 цвета с точностью до поворотов и отражений.
	Группа симметрий — диэдральная группа $D_9$, $|D_9| = 18$.
	Найдём мощности множеств неподвижных точек $|M^g|$ для каждого типа преобразований:
	\begin{itemize}
		\item $R_0$ (тождественное): 1 элемент.
		Все $2^9$ раскрасок остаются на месте. $|M^{R_0}| = 512$.
		\item $S_i$ (отражения): 9 элементов.
		Ось симметрии проходит через 1 вершину и середину противоположной стороны. Вершины разбиваются на 1 неподвижную и 4 пары симметричных.
		Для сохранения раскраски цвета в парах должны совпадать. Независимых выборов цвета: 5. $|M^{S_i}| = 2^5 = 32$.
		\item $R_i$ (повороты порядка 9): Это повороты на углы, кратные $2\pi/9$, взаимно простые с 9. Таких элементов $\varphi(9) = 6$.
		Орбита вершин одна, поэтому все вершины должны быть одного цвета. $|M^{R_i}| = 2^1 = 2$.
		\item $R_i$ (повороты порядка 3): Повороты на $120^\circ$ и $240^\circ$.
		Таких элементов 2. Вершины разбиваются на 3 независимых цикла по 3 вершины. $|M^{R_i}| = 2^3 = 8$.
	\end{itemize}
	Применим лемму Бёрнсайда:
	\[ N = \frac{1}{18} (1 \cdot 512 + 9 \cdot 32 + 6 \cdot 2 + 2 \cdot 8) = \frac{1}{18} (512 + 288 + 12 + 16) = \frac{828}{18} = 46 \] 
	Всего существует 46 различных ожерелий.
\end{Example}
